\documentclass[10pt,a4paper]{amsart}
\usepackage[margin=19mm]{geometry}
\usepackage{amsmath,amssymb,mathtools}
\usepackage[scr=rsfs,cal=euler]{mathalfa}
\usepackage{xfrac}
\usepackage[unicode,hidelinks,bookmarks=false]{hyperref}
\newcommand{\ie}{i.e.\ }
\newcommand{\cf}{cf.\ }
\newcommand{\resp}{respectively\ }
\newcommand{\Subsection}{\subsection}
\providecommand{\nobreakdash}{\nobreak\hbox{-}}
\setlength{\emergencystretch}{2em}
\multlinegap0pt
\usepackage{amssymb}
\usepackage{mathtools}
\usepackage{bm}
\usepackage{xfrac}
\usepackage[shortlabels]{enumitem}
\usepackage{float}
\newtheorem{proposition}{Proposition}
\title{Certified and accurate computation of function space norms of deep neural networks}
\author{Johannes Gr\"undler, Moritz Maibaum, Philipp Petersen}
\date{Source: arXiv:2603.06431v2 (CC BY 4.0)}
\begin{document}
\maketitle

\noindent\emph{Source and licence.} Certified and accurate computation of function space norms of deep neural networks. Version arXiv:2603.06431v2 (CC BY 4.0), distributed by arXiv under the Creative Commons Attribution 4.0 International licence, http://creativecommons.org/licenses/by/4.0/, as recorded in the arXiv licence metadata for this exact version and shown on the abstract page https://arxiv.org/abs/2603.06431v2. Full licence notice: \texttt{licenses/arxiv-2603.06431-CC-BY-4.0.txt}.\par\smallskip

\noindent\textit{Editorial scope.} Counted unit: the article's proposition hölder (source0.tex lines 1412--1414). The article's complete original proof of it is delivered with it (source0.tex lines 1415--1449, one \texttt{proof} environment of 35 lines). No mathematics is written, repaired or re-derived: the statement and the proof are the article's own lines, in the article's own notation, and no retained line is substituted or re-flowed.  Retained uncounted as the source-local context the proof needs: lines 490--492; lines 1187--1191; lines 1373--1376.  Nothing was omitted from any retained span, so the package carries no omission record.  No retained line contains a citation, so the wrapper carries no bibliography and no bibliography entry.  No retained line contains a figure environment, an image-inclusion command or drawing code, so no image is delivered and the package declares no asset dependency.  Editorial formatting only, in addition: an AMS \texttt{amsart} preamble that restates the article's own declarations and macros used by the retained lines, namely the environments \texttt{proposition} on separate counters, together with the standard packages those lines need.  The retained environments print in this document as Proposition 1 in order of appearance, whereas the article prints the same items with the numbering of its own full document.  The complete native source of the article as distributed in the arXiv e-print for this exact version is bundled unchanged as \texttt{sources/195872/source0.tex}.
		\begin{equation}\label{eq: matrix width inequality}
			w(AB) \leq ||A||w(B) + ||B||w(A)
		\end{equation}

	\begin{proposition}\label{prop: fval hölder}
		Let $\Phi : \mathbb{R}^{d_0} \to \mathbb{R}^{d_{L+1}}$ be a neural network with activation function $\sigma : \mathbb{R} \to \mathbb{R}$ and Hölder continuous interval enclosure $\Sigma:\mathbb{IR}\to \mathbb{IR}$ having exponent $\gamma_\Sigma \in (0,1]$ and constant $C_\Sigma > 0$. Then for $\ell=1,\dotsc,L+1$ the function $\mathrm{Fval}_{\Phi, \ell, \Sigma}$ is Hölder continuous on $\Omega$, i.e., there exist $C > 0$ and $\gamma \in (0,1]$ such that
		$$ w(\mathrm{Fval}_{\Phi, \ell, \Sigma}(K)) \leq C w(K)^{\gamma} $$
		for all $K \in \mathbb{IR}^{d_0}$ with $K \subset \Omega$.
	\end{proposition}

		\begin{equation}\label{eq: algo jac recursion}
			\mathrm{Jac}_{\Phi,\hat{\ell}, \Sigma, \Sigma'}(\hat{K}) = 
			\mathrm{Jac}_{\Phi,\hat{\ell}+1, \Sigma, \Sigma'}(\hat{K})  \operatorname{diag}(\Sigma'(\mathrm{Fval}_{\Phi , \hat{\ell}+1,\Sigma}(\hat{K}))) W^{(\hat{\ell})}
		\end{equation}

	\begin{proposition}\label{prop: jac hölder}
		Let $\Omega \in \mathbb{IR}^{d_0}$, $\Phi : \mathbb{R}^{d_0} \to \mathbb{R}^{d_{L+1}}$ be a neural network with differentiable activation function $\sigma : \mathbb{R} \to \mathbb{R}$ and Hölder continuous interval enclosures $\Sigma,\Sigma':\mathbb{IR}\to \mathbb{IR}$ of $\sigma,\sigma'$, respectively. Then for $\ell=0,\dotsc,L$ the function $\mathrm{Jac}_{\Phi,\ell, \Sigma, \Sigma'}$ is Hölder continuous on $\Omega$.
	\end{proposition}

	\begin{proof}
		Let $K\in \mathbb{IR}^{d_0}, \, K \subset \Omega$. We have $\mathrm{Jac}_{\Phi,L, \Sigma, \Sigma'}(K)=W^{(L)}$ and thus
		$$ w(\mathrm{Jac}_{\Phi,L, \Sigma, \Sigma'}(K)) = 0, $$
		which is in particular Hölder continuous for every exponent in $(0,1]$ and constant zero.
		Let $\ell \in \{0, \dotsc,L-1\}$ be arbitrary but fixed and assume
		\begin{equation}\label{eq: jac hölder induction hypothesis}
			w(\mathrm{Jac}_{\Phi,\ell+1, \Sigma, \Sigma'}(K)) \leq C_{\ell+1} w(K)^{\gamma_{\ell+1}}
		\end{equation}
		with $C_{\ell +1} >0$ and $\gamma_{\ell +1}\in (0,1]$.
		By Equation~\eqref{eq: algo jac recursion} we have
		\begin{align*}
			w(\mathrm{Jac}_{\Phi,\ell, \Sigma, \Sigma'}(K)) &=
			w\left( \mathrm{Jac}_{\Phi,{\ell}+1, \Sigma, \Sigma'}({K})  \operatorname{diag}(\Sigma'(\mathrm{Fval}_{\Phi , {\ell}+1, \Sigma}({K}))) W^{(\ell)}\right)
		\end{align*}
		and applying Equation~\eqref{eq: matrix width inequality} yields
		\begin{align*}
			w\left( \mathrm{Jac}_{\Phi,{\ell}+1, \Sigma, \Sigma'}({K})  \operatorname{diag}(\Sigma'(\mathrm{Fval}_{\Phi , {\ell}+1, \Sigma}({K}))) W^{(\ell)}\right) 
			&\leq ||\mathrm{Jac}_{\Phi,{\ell}+1, \Sigma, \Sigma'}({K})||w(\Delta^{(\ell+1)}(K)) \\ 
			&\quad + ||\Delta^{(\ell+1)}(K)|| w(\mathrm{Jac}_{\Phi,{\ell}+1, \Sigma, \Sigma'}({K})),
		\end{align*}
		where we define $\Delta^{(\ell+1)}(K) = \operatorname{diag}(\Sigma'(\mathrm{Fval}_{\Phi , {\ell}+1, \Sigma}({K}))) W^{(\ell)}$ for brevity.
		We have
		$$w(\Delta^{(\ell+1)}(K))  \leq ||\operatorname{diag}(\Sigma'(\mathrm{Fval}_{\Phi , {\ell}+1, \Sigma}({K})))|| w(W^{(\ell)}) + ||W^{(\ell)}||w(\operatorname{diag}(\Sigma'(\mathrm{Fval}_{\Phi , {\ell}+1, \Sigma}({K})))) 
		$$
		by Equation~\eqref{eq: matrix width inequality}. Considering $w(W^{(\ell)})=0$ and $w(\operatorname{diag}(\Sigma'(\mathrm{Fval}_{\Phi , {\ell}+1, \Sigma}({K})))) =w(\Sigma'(\mathrm{Fval}_{\Phi , {\ell}+1, \Sigma}({K})))$ we obtain
		\begin{align*}
			&||\operatorname{diag}(\Sigma'(\mathrm{Fval}_{\Phi , {\ell}+1, \Sigma}({K}))|| w(W^{(\ell)}) + ||W^{(\ell)}||w(\operatorname{diag}(\Sigma'(\mathrm{Fval}_{\Phi , {\ell}+1, \Sigma}({K})))) \\
			= &||W^{(\ell)}||w(\Sigma'(\mathrm{Fval}_{\Phi , {\ell}+1, \Sigma}({K}))).
		\end{align*}
		By assumption $\Sigma'$ is Hölder continuous and we showed in Proposition~\ref{prop: fval hölder} that $\mathrm{Fval}_{\Phi , {\ell}+1, \Sigma}({K})$ is Hölder continuous. Thus there exist $\tilde{C}_{\ell +1} > 0$ and $\tilde{\gamma}_{(\ell+1)} \in (0,1]$ independent of $K$ such that
		$$ ||W^{(\ell)}||w(\Sigma'(\mathrm{Fval}_{\Phi , {\ell}+1, \Sigma}({K}))) \leq \tilde{C}_{\ell +1}w(K)^{\tilde{\gamma}_{(\ell+1)}}. $$
		It remains to bound the matrix norms of $||\mathrm{Jac}_{\Phi,{\ell}+1, \Sigma, \Sigma'}({K})||$ and $||\Delta^{(\ell+1)}(K)||$. Since both $\mathrm{Jac}_{\Phi,{\ell}+1, \Sigma, \Sigma'}$ and $\Delta^{(\ell+1)}(K)$ are inclusion isotonic, we have
		$$||\mathrm{Jac}_{\Phi,{\ell}+1, \Sigma, \Sigma'}({K})|| \leq ||\mathrm{Jac}_{\Phi,{\ell}+1, \Sigma, \Sigma'}({\Omega})|| \quad \text{and} \quad ||\Delta^{(\ell+1)}(K)|| \leq ||\Delta^{(\ell+1)}(\Omega)||.$$
		Note, that the inclusion isotonicity of $\Delta^{(\ell+1)}$ follows from the inclusion isotonicity of $\mathrm{Fval}_{\Phi, \ell+1, \Sigma}$.
	\end{proof}

\end{document}
